\documentclass{article}
\usepackage{doi}

\usepackage{hyperref}
\usepackage{graphicx} % Required for inserting images
\usepackage{german}
\title{Das Proseminar Kombinatorik}
\author{Tuan Anh Nguyen}
\date{\today}

\begin{document}

\maketitle
\section{Themen}
Die vorzutragenden Themen basieren auf den Büchern \cite{aigner2006diskrete,hougardy2018algorithmische}:
\begin{enumerate}
    \item Dienstag, 14. April 2026 -- Jana Chung \\ Primzahltest und Berechenbarkeit (\cite[Kapitel 1]{hougardy2018algorithmische})
    \item Dienstag, 21. April 2026 -- Carl-Philip Deerberg\\
	Darstellungen ganze Zahlen (\cite[Kapitel 2]{hougardy2018algorithmische})
    \item Dienstag, 28. April 2026 -- Alisa Mirena\\Rechnen mit ganzen Zahlen (\cite[Kapitel 3]{hougardy2018algorithmische})
    \item Dienstag, 5. Mai 2026 -- Max Recker\\Approximative Darstellung reeller Zahlen (\cite[Kapitel 4]{hougardy2018algorithmische})
    \item Dienstag, 12. Mai 2026 -- Ayleen Peschke\\Rechnen mit Fehlern (\cite[Kapitel 5]{hougardy2018algorithmische})
    \item Dienstag, 19. Mai 2026 -- Sophie Hartmann\\Graphen (\cite[Kapitel 6]{hougardy2018algorithmische})
    \item Dienstag, 26. Mai 2026 -- Lennart Meier\\Einfache Graphenalgorithmen (\cite[Kapitel 7]{hougardy2018algorithmische})
    \item Dienstag, 2. Juni 2026 -- Maya Neudorf\\Sortieralgorithmen (\cite[Kapitel 8]{hougardy2018algorithmische})
    \item Dienstag, 9. Juni 2026 --  Kristin Ackermann\\Optimale Bäume und Wege ((\cite[Kapitel 9]{hougardy2018algorithmische}))    
   % \item Gauss-Elimination (\cite[Kapitel 11]{hougardy2018algorithmische})
    \item Dienstag, 16. Juni 2026 -- Aileen Hays \\Elementare Zählprinzipien und Permutationen (\cite[Abschnitte 1.1--1.3]{aigner2006diskrete})
    \item Dienstag, 23. Juni 2026 -- Valerie Höhle\\Rekursion und Existenzaussagen (\cite[Abschnitte 1.4--1.6]{aigner2006diskrete})
    \item Dienstag, 30. Juni 2026 -- Simon Boss\\Summation -- Teil 1 (\cite[Abschnitte 2.1--2.2]{aigner2006diskrete})
    \item Dienstag, 7. Juli 2026 -- Zsombor Hornyak\\Summation -- Teil 2 (\cite[Abschnitte 2.3--2.4]{aigner2006diskrete})
    \item Dienstag, 14. Juli 2026 -- Carla Holsen\\Erzeugende Funktionen (\cite[Kapitel 3]{aigner2006diskrete})
    \item Dienstag, 21. Juli 2026 -- Mira Sivulka\\ Gauss-Elimination (\cite[Kapitel 11]{hougardy2018algorithmische})
\end{enumerate}
\section{Voraussetzung} Um das Proseminar zu besuchen, braucht man nur die Schulwissen. Dennoch ist das Proseminar ohne mathematische Vorbildung anspruchsvoll.
\section{Anforderungen}
\begin{enumerate}
    \item aktive Teilnahme am Proseminar (mindestens 12 Termine)
    \item erfolgreicher Vortrag mit einer Ausarbeitung in \LaTeX
    \begin{itemize}
        \item Termin für die erste Version der Ausarbeitung: eine Woche vor dem Vortragtermin
        \item Termin für die letzte Version der Ausarbeitung: $8$ Wochen nach dem Vortragtermin
        \item Alle Versionen der Ausarbeitung (.tex, .bib, und .pdf files) müssen auf \url{www.github.com} hochgeladen   werden und zu den anderen TeilnehmerInnen zugreifbar sein.
    \end{itemize}
\end{enumerate}

\bibliographystyle{plain} % or alpha, unsrt, etc.
\bibliography{literature}
\end{document}
